\section{Finite information density and interferometric searches for Planckian noise}
\label{sec:holographic}

A simulation run on finite hardware can only represent a finite amount of
information about any region. Holographic entropy bounds say that, if
semiclassical gravity is a reliable guide, nature also stores only a finite
amount of information in a region, and that this amount scales with the area of
the boundary rather than with the volume. This section reviews what those bounds
imply about the ``resolution budget'' of any substrate that reproduces our
physics. It then turns to the one experimental programme that has tried to
detect a macroscopic consequence of finite spatial information: interferometric
searches for Planckian ``holographic noise''. The models tested in these
searches come from quantum gravity, not from assumptions about simulation. We
state this at the start because the popular description of these experiments as
tests of whether space is ``pixelated'' blurs the distinction.

\subsection{Holographic bounds as an information budget}

Bekenstein argued from the generalized second law (GSL) that a weakly
gravitating system of energy $E$ and radius $R$ in asymptotically flat space
obeys $S_{\rm matter}\le 2\pi E R$ in Planck units~\cite{Bekenstein1981universal};
we rely on Bousso's account of the argument, which also notes that whether the
GSL implies this bound ``remains controversial''~\cite{Bousso2002holographic}.
't~Hooft went further and argued that the observable degrees of freedom can be
described as Boolean variables on a two-dimensional lattice, about ``one Boolean
variable per Planckian surface element''~\cite{tHooft1993dimensional}.
Susskind recast this proposal as a ``screen'' of ``pixels'', each storing one
bit~\cite{Susskind1995world}. That is the origin of the pixel vocabulary, and it
was introduced as a heuristic, not as a claim about computation. In Bousso's
review the covariant entropy bound limits the information content of the
spacetime regions adjacent to a surface to $1.4\times10^{69}$ bits per square
metre, and the holographic principle is the assertion that this bound has its
origin in a quantum theory of gravity~\cite{Bousso2002holographic}. Bousso also
spells out the contrast with a naive lattice. One oscillator with $n$ states per
Planck volume gives $\sim n^{V}$ states, which scales with volume and is ``at
odds'' with the spherical entropy bound~\cite{Bousso2002holographic}.

These results bear on the simulation hypothesis in three ways. First, they give
a physically motivated upper bound on the memory a faithful simulation of a
region would need: a substrate that reproduces semiclassical gravity need not
store more than an area-law number of bits. For the whole observable universe,
Lloyd estimated at most $\sim10^{120}$ elementary operations on $\sim10^{90}$
bits ($\sim10^{120}$ bits if gravitational degrees of freedom are
included)~\cite{Lloyd2002computational}. Vazza used the holographic bound to
estimate $\sim3.5\times10^{124}$ bits and concluded that a Planck-resolution
simulation of the visible universe by a simulator sharing our physics is
impossible on energy grounds~\cite{Vazza2025astrophysical}. Vazza also notes
that a coarse-grained description of the cosmic web needs vastly fewer bits, so
the conclusion depends on the resolution assumed~\cite{Vazza2025astrophysical}.
The resource arguments themselves are reviewed in the section on computational limits.
% cross-link: Section 02 (computational resources); replace with \ref at merge.
Second, the bounds limit \emph{maximum entropy}. They do not imply that geometry
is discrete, or that a finite information density must leave a measurable
imprint. 't~Hooft himself remarked that the holographic ``blurring is small
compared to the uncertainties produced by the usual quantum mechanical
fluctuations''~\cite{tHooft1993dimensional}. Third, whether the information
limit could have \emph{macroscopic} observable consequences is the question that
motivated the experiments below.

\subsection{Hogan's holographic noise and the GEO600 episode}

Hogan proposed that a transverse indeterminacy of position, inferred from the
limits of measurement with Planck-wavelength radiation, would produce
``holographic noise''. This noise would have a transverse shear signature and a
flat strain power spectral density equal to the Planck time $t_P$, predicted
``with no parameters''~\cite{Hogan2008measurement}. The step from information to
noise is explicit. The normalization uses Nyquist--Shannon sampling with one
degree of freedom per $2l_P$ in each transverse direction on a light sheet, and
an experimental upper bound on the noise would therefore give ``a lower bound on
the number density of independent position eigenstates''~\cite{Hogan2008measurement}.
Hogan stated that the effect is ``not derived here from a fundamental theory''.
He compared the predicted spectrum with GEO600, not LIGO, arguing that LIGO's
arm cavities suppress the effect~\cite{Hogan2008measurement}, and he wrote that
the approximate agreement with ``otherwise unexplained noise in GEO600 motivates
further study''~\cite{Hogan2008measurement,Hogan2008indeterminacy}. A 2009
revision lowered the predicted amplitude by $\sqrt{\pi}$, to
$h=\sqrt{t_P/2\pi}=0.92\times10^{-22}\,{\rm Hz^{-1/2}}$ for GEO600. It also
withdrew an earlier claim about the spectral slope and predicted strong
correlation between nearly co-located interferometers~\cite{Hogan2009holographic}.

The proposal was contested. Smolyaninov argued that the derivation assumes
linear diffraction of Planck radiation, and that vacuum self-focusing would
reduce the angular uncertainty to $\sim\lambda_p/L$, ``19 orders of magnitude
smaller''~\cite{Smolyaninov2009level}. Amelino-Camelia called it ``a young
proposal still looking for some maturity''. He also stressed that it is unrelated
to earlier ``holography-inspired'' noise estimates of the spacetime-foam
type~\cite{AmelinoCamelia2013quantum}, so the phrase ``holographic noise'' names
two different proposals. On the experimental side, GEO600 papers of the period
report noise between roughly 100\,Hz and 1\,kHz that had not yet been
explained~\cite{Hild2009dcreadout,Luck2010upgrade}. The GEO600 team later
reported that in ``tuned DC readout'' experiments ``all of the observed noise
could be explained'' by known sources~\cite{Luck2010upgrade}; that paper does not
mention holographic noise. Amelino-Camelia concluded that ``the brief season of
the `GEO600 anomaly' \ldots\ is over''~\cite{AmelinoCamelia2013quantum}.
Hogan's revised analysis argued that folded arms (GEO600) and Fabry--Perot arms
(LIGO) do not amplify holographic noise, so that it was ``not currently ruled out
by either LIGO or GEO600''. He proposed a dedicated test: cross-correlating two
nearly co-located interferometers at frequencies near $c/2L$~\cite{Hogan2012interferometers}.
Separately, Kwon and Hogan showed that LIGO/Virgo stochastic-background limits
already exclude some \emph{metric}-fluctuation (spacetime-foam) models, whereas
their non-metric holographic model remained near the limits of the experiments
of the time~\cite{Kwon2016interferometric}.

Hogan's most explicitly informational statement is a non-refereed essay. It
opens ``If reality has finite information content, space has finite fidelity''
and posits a ``Planck broadcast'' bandwidth of about $2\times10^{43}$ bits per
second~\cite{Hogan2013now}. The essay uses computing metaphors (a ``cosmic
Internet Service Provider'', the ``operating system'' the universe ``runs on'').
It also rejects Planck-size pixels (``nature is not really pixelated in little
squares'') in favour of a transverse blur
$\langle x_\perp^2\rangle=(2.135\times10^{-18}\,{\rm m})^2(L/1\,{\rm m})$, far
larger than the Planck length~\cite{Hogan2013now}. It does not propose that the
universe is a simulation.

\subsection{The Fermilab Holometer}

The Holometer consisted of two co-located 39.06\,m power-recycled Michelson
interferometers (about 2\,kW on each beamsplitter, beamsplitters 0.91\,m apart,
no arm cavities). Their outputs were cross-correlated up to 25\,MHz to average
down uncorrelated shot noise~\cite{Chou2016first,Holometer2017holometer}.
Each interferometer had a shot-noise-limited sensitivity of
$2.1\times10^{-18}$\,m\,Hz$^{-1/2}$. After 145\,h of data the cross-spectrum
reached $2.1\times10^{-20}$\,m\,Hz$^{-1/2}$, surpassing strain power spectral
density $t_P$ for signal bandwidths above 11\,kHz~\cite{Chou2016first}.
The first analysis tested a ``speculative model'' of Planckian diffraction
from ``a putative fundamental Nyquist frequency $f_p=1/t_p$'' (the Kwon--Hogan
model~\cite{Kwon2016interferometric}). It excluded that model at $5.1\sigma$
statistical significance ($4.6\sigma$ including a 10\% calibration uncertainty),
equivalently a normalization below 44\% of the prediction at 95\% CL. The authors
stressed that this ``appl[ies] only to the spectral shape of the particular model''
and that the layout ``would not respond'' to rotational
correlations~\cite{Chou2016first}. The final straight-arm analysis (704\,h)
found no correlation. At $2\sigma$ it limited the amplitudes of a two-parameter
shear model to $|\beta_L|<0.10\,t_P$ and $|\beta_{2L}|<0.25\,t_P$, well below
the order-unity values predicted, and stated explicitly that it does not
constrain a newer Lorentz-invariant rotational model~\cite{Chou2017interferometric}.
A reconfiguration with one arm of each interferometer bent by $90^\circ$ tested
that rotational model~\cite{Hogan2017statistical} with 1098\,h of data. It
constrained its normalization to $\eta<0.25\,t_P$ at 95\% confidence for one
rotation axis~\cite{Richardson2021interferometric}. The authors note that the
planar geometry cannot test models whose uncertainty entangles all three spatial
directions, and that local quantum-field-theoretic models of quantum gravity
predict no effect in such a measurement~\cite{Richardson2021interferometric}.
The same data also gave the first direct limits on MHz stochastic gravitational
waves, which lie far above closure density~\cite{Holometer2017mhz}.

In short, the Holometer placed null bounds on specific models: two spectral
families of shear noise and one rotational model. It did not rule out Planckian
information limits in general. A Holometer collaborator later described the
first-generation design as resting on a ``naive premise'' that drew ``(justified)
criticism'' for not being Lorentz invariant, and reinterpreted its null result
as ``a verification of an exact symmetry''~\cite{Kwon2025phenomenology}. We could
not locate a published version of that Lorentz-invariance criticism.

\subsection{Geontropic fluctuations and next-generation experiments}

Verlinde and Zurek started from the covariant entropy bound and postulated
``Planck size pixels on the surface bounding a causally connected part of
space''. Energy fluctuations of these pixels would produce arm-length
fluctuations $\langle\delta L^2\rangle\sim l_pL$ with strong transverse
correlations~\cite{Verlinde2021observational}. Follow-up work tied the effect to
the area-law fluctuations of the modular Hamiltonian, computed in AdS/CFT and
modelled in flat space by a high-occupation bosonic field~\cite{Verlinde2020spacetime,Zurek2022vacuum,Verlinde2022modular}.
Zurek notes that ``it is not known how to precisely extend'' these ideas to flat
space~\cite{Zurek2022vacuum}. Li et al.\ recast published LIGO and Holometer data
as rough $3\sigma$ bounds on the normalization $\alpha$ (natural value
$\alpha\sim1$). They find $\alpha\lesssim3$ (LIGO) and $\alpha\lesssim0.7$
(Holometer) with an infrared cutoff, and $\alpha\lesssim0.1$ and
$\alpha\lesssim0.6$ without one. Contrary to Hogan's earlier assumption, they
also argue that these fluctuations do accumulate in Fabry--Perot
arms~\cite{Li2023interferometer}. Two new cross-correlation experiments target
this regime. GQuEST (Caltech/Fermilab) is a design for 5\,m interferometers with
photon-counting readout that evades the standard quantum limit. It projects
$\alpha<0.6$ at $3\sigma$ in 60\,h and $\alpha<0.1$ in 2160\,h at design
sensitivity~\cite{Vermeulen2025photon}. QUEST (Cardiff)~\cite{Vermeulen2021experiment}
has reported first limits on correlated length fluctuations between 13 and
80\,MHz from a $10^4$\,s run with 1.83\,m arms. It frames these as
stochastic-background bounds and quotes no constraint on a specific
quantum-gravity model~\cite{Patra2025broadband}. If the effect were real, it
would be a large background for Cosmic Explorer and the Einstein
Telescope~\cite{Bub2023quantum}. A tabletop demonstration has shown that
quantum-correlated light can outperform classical light in detecting a correlated
signal injected into two interferometers in the MHz band~\cite{Pradyumna2020twin}.

Other theorists dispute that an effect of this size is expected. Carney et al.\
find that the effective field theory of gravitons ``unambiguously predicts''
$\Delta L\sim\ell_{\rm pl}\sim10^{-35}$\,m. On this view a detection would signal
``a severe breakdown of effective quantum field theory''~\cite{Carney2026response}.
A recent preprint likewise finds graviton-induced noise spectra that are finite
but Planck-suppressed~\cite{Freidel2026geometric}. Reviews place these searches
among tests of particular phenomenological models of spacetime
fluctuations~\cite{Carlip2023spacetime,Sharmila2026signatures}.

\begin{table*}[t]
\centering
\footnotesize
\caption{Interferometric searches for Planckian, holographically motivated noise.
Sensitivities are as quoted by the cited papers (cross-spectrum values unless
stated). ``Proposal'' marks projections, not measurements.}
\label{tab:holographic-experiments}
\footnotesize\setlength{\tabcolsep}{3pt}
\begin{tabular}{p{2.1cm}p{2.5cm}p{2.9cm}p{2.7cm}p{3.6cm}p{1.6cm}}
\hline
Instrument & Observable & Sensitivity & Model tested & Result & Ref.\\
\hline
GEO600 (folded 600\,m arms) & DARM strain noise, $\sim$100--600\,Hz & peak $2.2\times10^{-22}$\,Hz$^{-1/2}$ at 530\,Hz & Hogan holographic noise (comparison, not a dedicated test) & Unexplained noise later attributed to known sources in DC readout; Hogan: not ruled out after arm-folding factor & \cite{Hogan2008measurement,Luck2010upgrade,Hogan2012interferometers}\\
LIGO/Virgo & stochastic-background cross-correlation (41.5--169.25\,Hz); published strain spectra & $\Omega_{\rm GW}<5.6\times10^{-6}$ (95\% CL) used for foam models & metric-fluctuation (foam) models; geontropic recast & ``random walk'' model excluded; white-noise model $a_0<0.06$; geontropic $\alpha\lesssim0.1$ (no IR cutoff) or $\lesssim3$ (IR cutoff) & \cite{Kwon2016interferometric,Li2023interferometer}\\
Holometer I (39\,m, straight arms, 2015--16) & DARM cross-spectrum up to 25\,MHz & $2.1\times10^{-20}$\,m\,Hz$^{-1/2}$ (145\,h) & Kwon--Hogan Planckian diffraction; shear-noise family & Excluded at $4.6\sigma$ ($<44\%$ at 95\% CL); $|\beta_L|<0.10\,t_P$, $|\beta_{2L}|<0.25\,t_P$ ($2\sigma$, 704\,h) & \cite{Chou2016first,Chou2017interferometric}\\
Holometer II (38.9\,m, bent arms) & rotational DARM cross-spectrum, 1.1--20\,MHz & $6.1\times10^{-21}$\,m\,Hz$^{-1/2}$; CSD$_{\delta h}<t_P/2$ (1098\,h) & Hogan--Kwon--Richardson rotational model & $\eta<0.25\,t_P$ (95\% CL, one axis) & \cite{Richardson2021interferometric}\\
QUEST (Cardiff, 1.83\,m) & DARM cross-spectrum, 13--80\,MHz & strain $3\times10^{-20}$\,Hz$^{-1/2}$ ($10^4$\,s) & none specific (stochastic length fluctuations) & Upper limits; no model constraint quoted & \cite{Patra2025broadband}\\
GQuEST (Caltech/Fermilab, 5\,m) & photon-counting readout, peak $\sim$15.6\,MHz & Proposal: nominal signal in several hours & Verlinde--Zurek geontropic fluctuations & Proposal: $\alpha<0.6$ in 60\,h; $\alpha<0.1$ in 2160\,h ($3\sigma$) & \cite{Vermeulen2025photon}\\
\hline
\end{tabular}
\end{table*}

\subsection{Relation to the simulation hypothesis}

The logical link is weak and one-directional. (i) A simulation with finite
memory must coarse-grain, but nothing requires the coarse-graining to show up as
stationary, correlated displacement noise at the MHz frequencies these
instruments probe. A simulator could just as well reproduce smooth semiclassical
geometry to any accuracy we can measure, so a null result says nothing about
simulation as such. (ii) The models that were actually tested (Hogan's shear and
rotational models, the Verlinde--Zurek fluctuations) are derived from holography
and quantum-gravity arguments, with normalizations fixed by entropy bounds. None
of them is derived from assumptions about a simulating computer. A positive
detection would therefore first be evidence about quantum gravity. (iii) None of
the theoretical or experimental papers reviewed here invokes the simulation
hypothesis~\cite{Hogan2008measurement,Hogan2012interferometers,Chou2016first,Chou2017interferometric,Richardson2021interferometric,Verlinde2021observational,Vermeulen2025photon,Patra2025broadband}.
The nearest approach is metaphor: the sampling-theorem language of the Holometer
analysis~\cite{Chou2016first} and Hogan's ``Planck broadcast''
essay~\cite{Hogan2013now}. The only explicit link we found is in a non-refereed
preprint. It suggests that if experiments such as the Holometer showed ``the
holographic universe'' to be true, this ``may be interpreted as an indication of
us participating in a simulation chain''~\cite{Neukart2022do}. The experiment it
cites~\cite{Chou2016first} reported a null result for one noise spectrum and
did not test holography in general. Moreover, holographic scaling, if confirmed,
would be at least as naturally read as a fact about quantum gravity as evidence
of an ``external programmer''. What these experiments do provide is an
operational upper bound on some forms of macroscopic, Planck-normalized position
noise, and the evidence reviewed here does not support treating them as tests of
the simulation hypothesis.
